\documentclass[11pt]{article}
% ===========================================================================
%  Shared preamble for all MPM2D worksheets — input right after \documentclass.
%  (Files beginning with "_" are skipped by mpm2d-worksheet-publish.mjs.)
% ===========================================================================
\usepackage[letterpaper,top=2.2cm,bottom=1.9cm,left=1.6cm,right=1.6cm,headheight=28pt,footskip=24pt]{geometry}
\usepackage{amsmath,amssymb}
\usepackage[T1]{fontenc}
\usepackage[utf8]{inputenc}
\usepackage{lmodern}
\usepackage{xcolor}
\usepackage[most]{tcolorbox}
\usepackage{enumitem}
\usepackage{multicol}
\usepackage{fancyhdr}
\usepackage{tikz}
\usetikzlibrary{calc}
\usepackage{pgfplots}
\pgfplotsset{compat=1.18}

\definecolor{exblue}{HTML}{4A90E2}
\definecolor{exbg}{HTML}{E6F3FF}
\definecolor{qorange}{HTML}{E69138}
\definecolor{qbg}{HTML}{FFF7CC}
\definecolor{keygray}{HTML}{475569}
\definecolor{keybg}{HTML}{F1F5F9}
\definecolor{iagreen}{HTML}{14653B}
\definecolor{titleblue}{HTML}{1D4ED8}
% Rotating palette so each worked example gets its own colour.
\definecolor{ex0}{HTML}{2563A0}\definecolor{ex1}{HTML}{3B7D3B}\definecolor{ex2}{HTML}{8A5A00}
\definecolor{ex3}{HTML}{A3327A}\definecolor{ex4}{HTML}{6D28A3}\definecolor{ex5}{HTML}{B08900}
\definecolor{ex6}{HTML}{0E7490}\definecolor{ex7}{HTML}{9A3412}\definecolor{ex8}{HTML}{15803D}

\pagestyle{fancy}
\fancyhf{}
\fancyhead[L]{\footnotesize NAME: \underline{\hspace{3.4cm}}}
\fancyhead[C]{\textbf{Integration Academy}\\[-2pt]{\scriptsize Math Simplified \textperiodcentered{} Success Amplified}}
\fancyhead[R]{\footnotesize MPM2D}
\fancyfoot[L]{\scriptsize dr.merie@IntegrationAcademy.ca}
\fancyfoot[C]{\scriptsize\thepage}
\fancyfoot[R]{\scriptsize IntegrationAcademy.ca}
\renewcommand{\headrulewidth}{1.2pt}
\renewcommand{\footrulewidth}{0.4pt}
\setlength{\parindent}{0pt}
\setlength{\parskip}{4pt}

% Examples are NOT breakable, so each example + its solution stays on one page.
\newtcolorbox{example}[2]{colframe=#1,colback=#1!7,coltitle=white,colbacktitle=#1,fonttitle=\bfseries,title={#2},arc=3pt,boxrule=1pt}
\newtcolorbox{qbox}[1]{colframe=qorange,colback=qbg,coltitle=white,colbacktitle=qorange,fonttitle=\bfseries,title={#1},arc=3pt,breakable,boxrule=1pt}
\newtcolorbox{keybox}{colframe=keygray,colback=keybg,coltitle=white,colbacktitle=keygray,fonttitle=\bfseries,title={$\checkmark$\ Answer Key},arc=3pt,breakable,boxrule=1pt}
\newtcolorbox{recapbox}{colframe=keygray,colback=keybg,coltitle=white,colbacktitle=keygray,fonttitle=\bfseries,title={Question Recap},arc=3pt,breakable,boxrule=1pt}
\newtcolorbox{ideabox}{colframe=titleblue,colback=titleblue!6,coltitle=white,colbacktitle=titleblue,fonttitle=\bfseries,title={The Big Ideas},arc=3pt,breakable,boxrule=1.2pt}
\newtcolorbox{learnbox}{colframe=iagreen,colback=iagreen!5,coltitle=white,colbacktitle=iagreen,fonttitle=\bfseries,title={Concept Essentials},arc=3pt,breakable,boxrule=1.2pt}
\newcommand{\lh}[1]{\par\smallskip{\bfseries\color{iagreen}#1}\par\nobreak\smallskip}

\newcommand{\soln}{\par\smallskip\textit{\textbf{Solution.}}\ }
\newcommand{\workspace}[1]{\par\vspace{3pt}\fcolorbox{black!30}{white}{\begin{minipage}[t][#1][t]{\dimexpr\linewidth-2\fboxsep\relax}\ \end{minipage}}\par}
\newcommand{\sech}[1]{\par\vspace{8pt}{\Large\bfseries\color{iagreen}#1}\par\nobreak\vspace{2pt}{\color{black!20}\hrule height1pt}\vspace{6pt}}

% A compact coordinate plot sized for an example box: \eplot{xmin}{xmax}{ymin}{ymax}{body}
\newcommand{\eplot}[5]{\begin{center}\begin{tikzpicture}
\begin{axis}[width=6.8cm,height=4.4cm,axis lines=middle,xlabel={\tiny$x$},ylabel={\tiny$y$},
grid=both,grid style={gray!16},xmin=#1,xmax=#2,ymin=#3,ymax=#4,
xtick distance=2,ytick distance=2,tick label style={font=\tiny},axis line style={-},samples=120]
#5
\end{axis}\end{tikzpicture}\end{center}}

% A sine-curve plot (degrees on x, 0..360): \splot{ymin}{ymax}{body}
\newcommand{\splot}[3]{\begin{center}\begin{tikzpicture}
\begin{axis}[width=8.6cm,height=4.6cm,axis lines=middle,xlabel={\tiny$x\,(^\circ)$},ylabel={\tiny$y$},
grid=both,grid style={gray!16},xmin=0,xmax=360,ymin=#1,ymax=#2,
xtick distance=90,ytick distance=1,tick label style={font=\tiny},axis line style={-},samples=180]
#3
\end{axis}\end{tikzpicture}\end{center}}

% A small coordinate plot: \plot{xmin}{xmax}{ymin}{ymax}{body}
\newcommand{\plot}[5]{\begin{center}\begin{tikzpicture}
\begin{axis}[width=8.4cm,height=6cm,axis lines=middle,xlabel={\scriptsize$x$},ylabel={\scriptsize$y$},
grid=both,grid style={gray!16},xmin=#1,xmax=#2,ymin=#3,ymax=#4,
xtick distance=2,ytick distance=2,tick label style={font=\tiny},axis line style={-}]
#5
\end{axis}\end{tikzpicture}\end{center}}

% Right triangle (angle theta at bottom-left, right angle at bottom-right):
%   \rtri{drawBase}{drawHeight}{oppLabel}{adjLabel}{hypLabel}
\newcommand{\rtri}[5]{\begin{center}\begin{tikzpicture}[scale=0.85]
\coordinate (A) at (0,0);\coordinate (B) at (#1,0);\coordinate (C) at (#1,#2);
\draw[thick,exblue] (A)--(B)--(C)--cycle;
\draw[black] ($(B)+(-0.32,0)$)--($(B)+(-0.32,0.32)$)--($(B)+(0,0.32)$);
\node[below] at ($(A)!0.5!(B)$){#4};
\node[right] at ($(B)!0.5!(C)$){#3};
\node[above left] at ($(A)!0.5!(C)$){#5};
\node at (0.7,0.22){$\theta$};
\end{tikzpicture}\end{center}}

% General (oblique) triangle with vertex labels and opposite-side labels:
%   \gtri{Alabel}{Blabel}{Clabel}{aSide}{bSide}{cSide}
\newcommand{\gtri}[6]{\begin{center}\begin{tikzpicture}[scale=0.85]
\coordinate (A) at (0,0);\coordinate (B) at (5,0);\coordinate (C) at (1.6,3);
\draw[thick,exblue] (A)--(B)--(C)--cycle;
\node[below left] at (A){#1};\node[below right] at (B){#2};\node[above] at (C){#3};
\node[below] at ($(A)!0.5!(B)$){#6};
\node[right] at ($(B)!0.5!(C)$){#4};
\node[left] at ($(A)!0.5!(C)$){#5};
\end{tikzpicture}\end{center}}

\fancyhead[R]{\footnotesize MCR3U \textperiodcentered{} 1.2}
\begin{document}

\begin{center}
{\LARGE\bfseries\color{titleblue}Worksheet 1.2 --- Domain and Range}\\[2pt]
{\small Grade 11 Functions, University (MCR3U) \textperiodcentered{} Unit 1: Characteristics of Functions}
\end{center}

\begin{ideabox}
The domain is all allowed inputs $x$; the range is all outputs $y$.
\begin{itemize}[leftmargin=*,itemsep=2pt]
  \item Denominators cannot be zero.
  \item Square roots need a non-negative radicand.
  \item A parabola's range starts at its vertex.
\end{itemize}
\end{ideabox}

\begin{learnbox}
\lh{What they mean}The \textbf{domain} is the set of allowed inputs ($x$-values); the \textbf{range} is the set of resulting outputs ($y$-values).
\lh{Finding restrictions}Exclude inputs that break the rule: values making a denominator $0$, or a negative under a square root, are not in the domain.
\lh{Reading from a graph}Scan left-to-right for the domain and bottom-to-top for the range, watching for arrows (continuing) versus open or closed endpoints.
\end{learnbox}

\sech{Worked Examples}

\begin{example}{ex0}{Example 1 --- Square-root domain}
Domain of $\sqrt{x-2}$?\soln The radicand must be $\ge0$: $x-2\ge0\Rightarrow x\ge2$. The curve starts at $(2,0)$:\eplot{-1}{8}{-1}{4}{\addplot[exblue,very thick,domain=2:8]{sqrt(x-2)};\addplot[red,only marks,mark=*,mark size=1.6pt] coordinates {(2,0)};}
\end{example}

\begin{example}{ex1}{Example 2 --- Square-root range}
Range of $\sqrt{x-2}$?\soln A square root is never negative, so outputs are $y\ge0$.
\end{example}

\begin{example}{ex2}{Example 3 --- Quadratic range}
Range of $x^2+1$?\soln The vertex $(0,1)$ is the lowest point, so $y\ge1$:\eplot{-3}{3}{-1}{8}{\addplot[exblue,very thick,domain=-2.6:2.6]{x^2+1};\addplot[red,only marks,mark=*,mark size=1.6pt] coordinates {(0,1)};}
\end{example}

\begin{example}{ex3}{Example 4 --- Rational domain}
Domain of $\dfrac{1}{x-3}$?\soln The denominator cannot be zero: $x-3\ne0\Rightarrow x\ne3$.
\end{example}

\begin{example}{ex4}{Example 5 --- Polynomial domain}
Domain of $x^2-5$?\soln A polynomial is defined everywhere: all real numbers.
\end{example}

\begin{example}{ex5}{Example 6 --- Reflected range}
Range of $-x^2$?\soln It opens down with vertex $(0,0)$ as the highest point: $y\le0$.
\end{example}

\begin{example}{ex6}{Example 7 --- Domain}
Domain of $\sqrt{x+5}$?\soln $x+5\ge0\Rightarrow x\ge-5$.
\end{example}

\begin{example}{ex7}{Example 8 --- Rational domain}
Domain of $\dfrac{1}{x+2}$?\soln $x+2\ne0\Rightarrow x\ne-2$.
\end{example}

\begin{example}{ex8}{Example 9 --- Range}
Range of $x^2-3$?\soln Vertex $(0,-3)$ is the minimum, so $y\ge-3$.
\end{example}

\clearpage
\sech{Your Turn --- Show Your Work}
For each question, show all steps clearly in the space provided.

\begin{qbox}{Question 1}
Domain of $\sqrt{x+5}$?\workspace{2.4cm}
\end{qbox}

\begin{qbox}{Question 2}
Domain of $\dfrac{1}{x+2}$?\workspace{2.4cm}
\end{qbox}

\begin{qbox}{Question 3}
Range of $x^2-3$?\workspace{2.4cm}
\end{qbox}

\begin{qbox}{Question 4}
Range of $-x^2$?\workspace{2.4cm}
\end{qbox}

\begin{qbox}{Question 5}
Domain of $\sqrt{2x-6}$?\workspace{2.4cm}
\end{qbox}

\begin{qbox}{Question 6}
Domain of $x^2+1$?\workspace{2.4cm}
\end{qbox}

\begin{qbox}{Question 7}
Range of $x^2+4$?\workspace{2.4cm}
\end{qbox}

\begin{qbox}{Question 8}
Domain of $\dfrac{1}{x-1}$?\workspace{2.4cm}
\end{qbox}

\begin{qbox}{Question 9}
Domain of $\sqrt{x-7}$?\workspace{2.4cm}
\end{qbox}

\begin{qbox}{Question 10}
Range of $-x^2+2$?\workspace{2.4cm}
\end{qbox}

\begin{qbox}{Question 11}
Domain of $3x-1$?\workspace{2.4cm}
\end{qbox}

\begin{qbox}{Question 12}
Range of $\sqrt{x}$?\workspace{2.4cm}
\end{qbox}

\begin{qbox}{Question 13 --- Challenge}
State the domain and range of $\sqrt{x-1}+2$.\workspace{2.6cm}
\end{qbox}

\clearpage
\sech{Question Recap \& Answer Key}

\begin{recapbox}
\begin{multicols}{2}
\small
\begin{enumerate}[leftmargin=*,itemsep=3pt]
  \item Domain of $\sqrt{x+5}$?
  \item Domain of $\dfrac{1}{x+2}$?
  \item Range of $x^2-3$?
  \item Range of $-x^2$?
  \item Domain of $\sqrt{2x-6}$?
  \item Domain of $x^2+1$?
  \item Range of $x^2+4$?
  \item Domain of $\dfrac{1}{x-1}$?
  \item Domain of $\sqrt{x-7}$?
  \item Range of $-x^2+2$?
  \item Domain of $3x-1$?
  \item Range of $\sqrt{x}$?
  \item State the domain and range of $\sqrt{x-1}+2$.
\end{enumerate}
\end{multicols}
\end{recapbox}

\begin{keybox}
\begin{multicols}{2}
\small
\begin{enumerate}[leftmargin=*,itemsep=3pt]
  \item $x\ge-5$
  \item $x\ne-2$
  \item $y\ge-3$
  \item $y\le0$
  \item $x\ge3$
  \item all reals
  \item $y\ge4$
  \item $x\ne1$
  \item $x\ge7$
  \item $y\le2$
  \item all reals
  \item $y\ge0$
  \item domain $x\ge1$, range $y\ge2$
\end{enumerate}
\end{multicols}
\end{keybox}

\end{document}
